% Generated draft: requires mathematical and source review.
\documentclass[11pt]{article}
\usepackage[margin=25mm]{geometry}
\usepackage{fontspec}
\setmainfont{Latin Modern Roman}
\usepackage{amsmath,amssymb,mathtools,unicode-math}
\setmathfont{Latin Modern Math}
\usepackage{enumitem,needspace,xurl,hyperref}
\setlength{\parindent}{0pt}
\setlength{\parskip}{5pt}
\setlength{\emergencystretch}{3em}
\newcommand{\N}{\mathbb{N}}
\newcommand{\Z}{\mathbb{Z}}
\newcommand{\Q}{\mathbb{Q}}
\newcommand{\R}{\mathbb{R}}
\newcommand{\C}{\mathbb{C}}
\newcommand{\F}{\mathbb{F}}
\newcommand{\A}{\mathbb{A}}
\newcommand{\PP}{\mathbb{P}}
\newcommand{\E}{\mathbb{E}}
\newcommand{\eps}{\varepsilon}
\newcommand{\dd}{\,\mathrm{d}}
\begin{document}
\section*{UnsolvedMath ID: 800487}
\textbf{Source problem code:} catalogue entry
\subsection{Definitions and mathematical statement}
The source asks for a characterization of optical local symplectic transformations of the space of oriented lines. The following makes the standard Euclidean interpretation explicit; the supplied record does not specify the ambient metric, the range of dimensions, or the precise convention for local realizability.

Let the ambient space be Euclidean \(\R^n\), with \(n\geq 2\). Its space of oriented affine lines is
\[
\mathcal L_n=\bigl(\R^n\times S^{n-1}\bigr)/\sim,
\qquad (q,v)\sim(q+tv,v)\quad(t\in\R).
\]
Here \(v\) is the unit direction of the line. Identify tangent and cotangent vectors using the Euclidean metric. On \(T^*\R^n\), put \(H(q,p)=\tfrac12\lVert p\rVert^2\). If \(\iota:H^{-1}(1/2)\hookrightarrow T^*\R^n\) is inclusion and \(\pi:H^{-1}(1/2)\to\mathcal L_n\) is the quotient by the Hamiltonian flow of \(H\), the reduced symplectic form \(\omega\) is characterized by
\[
\pi^*\omega=\iota^*\omega_{\mathrm{can}}.
\]
Thus \(\dim\mathcal L_n=2n-2\).

A smooth mirror is a smooth hypersurface \(M\subset\R^n\). Choose an open family of oriented lines and a smooth branch assigning to each line a transverse intersection point \(x\in M\). If \(v\) is the incoming unit direction and \(\nu(x)\) is either unit normal to \(M\), reflection sends the line through \(x\) with direction \(v\) to the line through \(x\) with direction
\[
v'=v-2\langle v,\nu(x)\rangle\nu(x).
\]
The formula is independent of the sign of \(\nu\). On a suitable open domain, this defines a local symplectomorphism \(R_M\) of \((\mathcal L_n,\omega)\).

Under the interpretation that optical realizability means finite composition of such local reflection maps, call a local symplectomorphism optical when it is representable as
\[
F=R_{M_k}\circ\cdots\circ R_{M_1}
\]
for finitely many smooth mirrors, with chosen intersection branches and compatible open domains throughout the composition.

\textbf{Problem (Serge Tabachnikov, Problem 9).} Characterize the optical transformations within the local symplectomorphisms of \((\mathcal L_n,\omega)\): give necessary and sufficient conditions for a given local symplectomorphism to admit such a finite mirror representation.

The source does not settle whether equality is required on the entire specified domain or only after restriction to a neighborhood, nor whether a representation must satisfy additional physical itinerary conditions, such as forward travel between successive mirrors and avoidance of unintended collisions. These conventions remain unresolved and can affect the class being characterized.

The source motivates the question by comparing the \(n-1\) variables describing a mirror locally with the \(2n-2\) variables of a Hamiltonian generating an infinitesimal symplectic transformation. Its expectation that the optical class is very small is heuristic; no precise notion of smallness is supplied, and no such assertion is included in the problem statement.
\subsection{Short English statement}
Characterize the local symplectic maps of the space of oriented lines that can be realized by finitely many smooth mirror reflections. The displayed formulation uses the standard Euclidean model. The source leaves the ambient setting, local domain convention, and physical itinerary requirements insufficiently specified.
\subsection{Sources}
\begin{itemize}
\item Local UnsolvedMath registry, permanent ID \(800487\), source code \texttt{catalogue entry}; supplied title: \emph{Billiards and Quantitative Symplectic Geometry --- Serge Tabachnikov --- Problem 9}. The registry labels the problem \emph{open}; this is a supplied status, not an independently verified mathematical or literature assessment.
\item Albach, Barraud et al., \emph{Open Problems in Billiards and Quantitative Symplectic Geometry} (2026), Serge Tabachnikov, Problem 9. Supplied source URL: \url{https://arxiv.org/abs/2602.12896}. The record attributes the source list to Albach, Barraud et al. and identifies Serge Tabachnikov as contributor. The supplied provenance records access on 2026-10-16 and extraction by \texttt{arxiv\_latex}; these are supplied provenance claims, not evidence of independent access or verification for this draft.
\item The supplied statement refers to results under the labels \texttt{[PTT,Pe,Ta20,Gl1]}. No corresponding bibliographic entries or links are supplied, so these references cannot be identified reliably from this record alone.
\end{itemize}
\end{document}
